%% abstract.tex
\begin{abstract}
Este proyecto propone la aplicación de técnicas de aprendizaje automático para el análisis predictivo del riesgo de corrupción en la contratación pública en Colombia. A partir de datos de SECOP II, se seleccionarán variables relevantes para entrenar modelos de clasificación que permitan identificar contratos con alta probabilidad de presentar irregularidades. La metodología incluye validación cruzada y evaluación mediante métricas como AUC-ROC y F1-score. El propósito es desarrollar un sistema de alertas tempranas que fortalezca la vigilancia preventiva y contribuya a una gestión más transparente y eficiente de los recursos públicos.
\end{abstract}

\begin{IEEEkeywords}
Elefantes Blancos, Sistema General de Regalías, Contratación pública, SECOP II, 
Aprendizaje automático, XGBoost, Random Forest, Alertas tempranas, 
Detección de anomalías, Auditoría preventiva
\end{IEEEkeywords}